\chapter{Protocolo de red v1}\label{cap:protocolo}

Resumen del protocolo; la especificación completa está en
\cmd{PROTOCOL-v1.md}, incluida en cada release. Cualquier cliente que la
implemente funciona con el dispositivo, en cualquier lenguaje.

\section{Transporte}

HTTP/1.1 sobre TCP, puerto 8080 (8443 con TLS). Control y estado en JSON; los
cuadros en un stream binario propio. Todas las rutas llevan el prefijo
\cmd{/v1}: dentro de v1 sólo se agregan campos y endpoints, nunca se quitan.

\section{Endpoints}

\begin{center}
\begin{tabularx}{\linewidth}{@{}l X@{}}
\toprule
\ep{GET}{/v1/health} & disponible sin token; \cmd{\{"ok": true\}} \\
\ep{GET}{/v1/device} & versiones (servicio, driver, bitstream, núcleo), capacidades, límites, TLS \\
\ep{GET}{/v1/status} & estado (\cmd{idle}, \cmd{decoding}, \cmd{error}), decode en curso, última secuencia \\
\ep{POST}{/v1/decode} & el clip en el cuerpo; la respuesta es el stream de cuadros \\
\ep{POST}{/v1/reset} & recuperación: resetea el núcleo y aborta el decode en curso \\
\ep{*}{/v1/debug/...} & diagnóstico, sólo con \cmd{-{}-debug}, sin garantía de compatibilidad \\
\bottomrule
\end{tabularx}
\end{center}

Un decode a la vez: un segundo \ep{POST}{/v1/decode} recibe 409 mientras el
primero corre. Los demás endpoints responden siempre.

\section{El stream de cuadros}

Una sucesión de registros, enteros \emph{little-endian}:

\begin{center}
\begin{tikzpicture}[x=11.5mm, font=\sffamily\scriptsize,
    f/.style={draw=#1, fill=#1!10, minimum height=9mm, anchor=west, inner sep=1pt, align=center}]
  % row 1: the 12-byte header every record starts with, one unit per byte
  \foreach \x/\w/\t in {0/4/{magic\\\texttt{"M2FR"}}, 4/1/tipo, 5/1/flags,
                         6/2/{header\\\_len}, 8/4/{payload\_len}}
    \node[f=m2blue, minimum width=\w*11.5mm] at (\x, 0) {\t};
  \foreach \b in {0,4,5,6,8,12} \node[lbl, below] at (\b, -0.48) {\b};
  \draw[decorate, decoration={brace, amplitude=4pt}, m2grey]
    (0, 0.55) -- node[above=4pt, lbl]{header común: 12 bytes} (12, 0.55);
  % row 2: what follows
  \node[f=m2teal, minimum width=5*11.5mm] at (0, -1.7)
    {campos del tipo\\FRAME: 20 bytes · END: ninguno};
  \node[f=m2amber, minimum width=7*11.5mm] at (5, -1.7)
    {payload\\FRAME: cuadro I420 · END: resumen JSON};
  \draw[decorate, decoration={brace, mirror, amplitude=4pt}, m2grey]
    (0, -2.25) -- node[below=4pt, lbl]{\texttt{header\_len} $-$ 12} (5, -2.25);
  \draw[decorate, decoration={brace, mirror, amplitude=4pt}, m2grey]
    (5, -2.25) -- node[below=4pt, lbl]{\texttt{payload\_len}} (12, -2.25);
\end{tikzpicture}
\end{center}

\begin{itemize}
  \item \textbf{FRAME} (tipo 1): índice de presentación, índice de decodificación,
    ancho, alto, tipo de imagen y estructura; el payload es el cuadro en I420.
  \item \textbf{END} (tipo 2): último registro, con un resumen en JSON:
    cuadros enviados, si el decode quedó completo, bits de error del
    decodificador y tiempos por etapa.
\end{itemize}

Un cliente debe usar \cmd{header\_len} y \cmd{payload\_len} y saltear los
campos y tipos de registro que no conozca: así crece el protocolo sin romper
clientes viejos.

\section{Códigos de error}

\begin{center}
\begin{tabular}{@{}c l l@{}}
\toprule
\textbf{HTTP} & \textbf{código} & \textbf{cuándo} \\
\midrule
400 & \cmd{bad\_request} & parámetro inválido, o el cuerpo no es MPEG-2 \\
401 & \cmd{unauthorized} & token ausente o incorrecto \\
404 & \cmd{not\_found} & ruta desconocida \\
409 & \cmd{busy} & ya hay un decode en curso \\
429 & \cmd{too\_many\_attempts} & demasiados tokens incorrectos \\
503 & \cmd{decoder\_unavailable} & el hardware no responde \\
\bottomrule
\end{tabular}
\end{center}
